\section{CrS-Aware Aggregation}\label{sec:agg}
% \abol{update}
% \paragraph{Coordination Mechanisms}.

As illustrated in Figure~\ref{fig:sys-arch}, after peer interactions, each agent $a_i\in A_t$ generates an output.
Existing coordination mechanisms (discussed in Appendix~\ref{app:coord}) integrate these outputs into an answer to the user query, following the strategies such as majority voting.

Building on top of the existing aggregation schemes, our system adds {\em credibility-scores} (CrS) to make the final output more reliable and robust against adversaries and low-performing agents.

Formally, the credibility score $CrS^{(j)}\in [0,1]$
of an agent $a_j\in \mathcal{A}$ is a {non-negative} number that reflects how reliable the system views the agent $a_i$ according to its performance in the previous query answering rounds. 

The credibility scores can be used in various coordination mechanism by replacing the unweighted aggregation with the weighted aggregation using the CrS scores.
Without loss of generality, in the following, we illustrate their integration into two integration mechanisms:

% \sana{Why are we talking about majority voting here? We do not have experiments on the integration of CrS into Majority Voting. So this should be deleted.}

% \vspace{1mm}
% \noindent{\em (a) majority vote:}
% In settings such as selecting between a set of options, majority voting {\em counts} the number of votes each option receives, and returns the answer with the maximum number of votes. Using the credibility scores, the answer with the maximum sum of credibility scores is returned. That is

% \vspace{-7mm}
% \[o_t = \text{argmax}_{o'\in \mathcal{O}(q_t)} \sum_{i : O(a_i,q_t)=o'} CrS^{(i)}\]

% \vspace{-2mm}\noindent
% Where $\mathcal{O}(q_t)$ is the universe of possible outputs for $q_t$, and $O(a_i,q_t)$ is the output of agent $a_i$ for $q_t$. %\mohsen{Is this CrS instead of SrC?}

\vspace{1mm}
\noindent{\em (a) centroid-based aggregation:} \label{sec:agg-requal}
% For settings where the agents' outputs are in the form of natural language and unstructured,
\cite{requal} proposes an aggregation strategy that first finds the centroid of the generated outputs in the embedding space, and then returns the closest answer to the centroid as the final output (see Appendix~\ref{app:related} for more details).
We use the CrS scores to find the {\em CrS-aware centroid} $\vec{x}^+$ as the weighted average of the generated outputs:

\vspace{-5mm}
\begin{equation}
\vec{x}^{+}=\frac{1}{N} \sum_{i : a_i\in A_t} {CrS}_{t-1}^{(i)} \vec{v}\left( O(a_i,q_t)\right)
\end{equation}

\vspace{-2mm}\noindent 
Where $O(a_i,q_t)$ is the output of agent $a_i$ for $q_t$, $\vec{v}(o)$ is the embedding of an output $o$, and ${CrS}_{t-1}^{(i)}$ is the current credibility score of $a_i$. 

\vspace{1mm}
\noindent{\em (b) LLM-assisted aggregation:}
Instead of using a specific formula for aggregation, one can use an LLM for this step, where in addition to the outputs $o_1,\cdots,o_N$, the {\em CrS scores} of the participating agents are sent to a {\bf Credibility‑aware Coordinator LLM}, which we trust.
The coordinator then aggregates the individual outputs and generates the final output while considering the CrS scores.

% \vspace{1mm}
% \noindent{\em (c) LLM-assisted aggregation:}
% The centroid-based aggregation does not consider the {\bf influence} of other agents in the output generated by an agent $a_i$. Therefore, while effective for settings such as no-connection topologies, it fails to consider adversary effects in individual outputs.  \sana{This part does not seem correct to me.}
% LLM-assisted aggregation can resolve this issue, where in addition to the outputs $o_1,\cdots,o_N$, the {\em dialogue log} alongside the {\em CrS scores} of the participating agents are sent to a {\bf Credibility‑aware Coordinator LLM}, which we trust.
% The coordinator then aggregates the individual outputs and generates the final output while considering the CrS scores and the dialogue log. \sana{We do not send the dialogue log to the coordinator; it is sent to the judge to compute the Contribution Scores. The coordinator only receives the final answers along with the agents’ credibility scores.}
